\documentclass{article}
\usepackage{dirtreex}
\makeatletter
\let\audit@preprocess\dte@preprocess
\def\dte@preprocess{%
  \typeout{AUDIT PREPROCESS tree=\the\dte@tnum; count=\the\dte@cnt; entry1=\dte@eget{1}{n}}%
  \ifnum\dte@cnt>1 \typeout{AUDIT entry2=\dte@eget{2}{n}}\fi
  \audit@preprocess}
\newlength\AuditLength
\def\AuditLiteralDimension{0pt}
\setlength\AuditLength{\AuditLiteralDimension}
\typeout{AUDIT SETLENGTH=\the\AuditLength}
\begin{document}
\begin{dirtreex}[pagebreak=false,box=false]
  \file{OUTER-BEFORE}{}
  \begin{dirtreex}[pagebreak=false,box=false]
    \file{INNER}{}
  \end{dirtreex}
  \file{OUTER-AFTER}{}
\end{dirtreex}
\typeout{AUDIT POST-NEST tree=\the\dte@tnum; count=\the\dte@cnt; entry1=\dte@eget{1}{n}}
\renewcommand\DirtreexFormatName[1]{%
  \ifnum#1=1
    \typeout{AUDIT FORMAT capture-flag=\ifdte@inenv TRUE\else FALSE\fi}%
    \file{ADDED-DURING-RENDER}{}%
  \fi
  \DirtreexGetField{#1}{n}}
\begin{dirtreex}[pagebreak=false,box=false]
  \file{ORIGINAL}{}
\end{dirtreex}
\typeout{AUDIT POST-FORMAT count=\the\dte@cnt; entry2=\dte@eget{2}{n}}
\end{document}
